\documentclass[sigconf]{acmart}
% Competition template settings. Source template:
% https://www.overleaf.com/read/vjztpvrnbrdh#43fbbe
\setcopyright{none}
\settopmatter{printacmref=false}
\acmConference[Kaggle Competition]{Filament Segmentation Challenge}{2026}{\href{https://www.kaggle.com/competitions/filament-segmentation-2026}{kaggle.com/competitions/filament-segmentation-2026}}
\acmBooktitle{Competition Report}
\acmDOI{}
\acmISBN{}
\usepackage{booktabs}
\usepackage{balance}
\usepackage{hyperref}
% Preserve the official geometry while avoiding text crossing narrow columns.
\emergencystretch=1.5em


\begin{document}
\title{Robust Radial Normalization for Solar Filament Instance Segmentation}
\subtitle{A Solution to the Solar Filament Segmentation Challenge 2026}
\author{ILoveBuns}
\affiliation{\country{China}}
\email{1280401315+ILoveBuns@users.noreply.github.com}

\begin{abstract}
We present a deterministic, CPU-friendly instance-segmentation pipeline for
solar filaments. The method removes center-to-limb brightness variation with
robust radial normalization, identifies dark filament candidates relative to
the visible solar disk, separates them into connected components, filters small
components, and emits compressed COCO run-length encodings. Our strongest
official Kaggle submission is platform reference \texttt{55113085}, with a
public score of \textbf{0.48}. This is higher than our later Mask R-CNN (0.31)
and frozen public YOLO plus U-Net (0.34) experiments. The classical pipeline is
therefore the selected final method. Its principal advantages are
reproducibility, low compute cost, and transparent failure modes.
\end{abstract}
\maketitle

\section{Introduction}
This report describes our submitted solution to the Solar Filament
Segmentation Challenge 2026 \cite{competition}, organized by the Earth-Space AI
Research Lab. The challenge uses GONG H-$\alpha$ observations and evaluates
pixel-level masks for individual solar filaments with Panoptic Quality (PQ)
\cite{kirillov2019panoptic}. In addition to overlap, the evaluation penalizes
false detections, missed instances, fragmentation, and over-merging.

Solar full-disk images exhibit strong center-to-limb brightness variation. A
fixed global threshold therefore confounds the dark limb and background with
actual filament material. Our contribution is an intentionally simple,
inspectable correction and grouping pipeline. We do not claim a new neural
architecture. Instead, we report the method that produced our highest verified
platform score, distinguish official scores from local checks, and provide an
end-to-end notebook that needs no weights or private supporting files.

\section{Methodology}\label{sec:methodology}
\subsection{Input and radial normalization}
Official JPEG inputs are converted to grayscale. For an image of height $H$ and
width $W$, we use its geometric center and a disk radius
$r_d=0.48\min(H,W)$. Disk pixels are divided into 96 radial bins. For every bin
$b$, the background profile $P_b$ is its upper-quartile intensity. The upper
quartile is less sensitive than the mean to dark filament contamination. A
pixel $x$ in bin $b(x)$ is normalized as
\begin{equation}
  N(x)=\frac{I(x)}{\max(P_{b(x)},10^{-6})}.
\end{equation}
This operation removes broad radial illumination variation while preserving
localized dark structures. Pixels outside the estimated disk are never
considered filament candidates.

\subsection{Candidate extraction and instance generation}
Within the disk, we compute the 25th percentile $q_{0.25}$ of normalized
intensity and set
\begin{equation}
  \tau=\min(q_{0.25},0.95).
\end{equation}
Pixels satisfying $N(x)\leq\tau$ form the candidate foreground. Four-connected
component labeling converts this foreground into instances. During submitted
inference, components smaller than 20 pixels are removed and no more than the
64 largest components per image are retained. Accepted masks are encoded as
compressed COCO RLE \cite{lin2014coco}. Sorting input paths and component sizes
makes the output deterministic for fixed data and dependencies.

The complete prediction path is implemented in \texttt{solarfil/infer.py} and
\texttt{solarfil/segment.py}. The supplied final notebook installs the pinned
environment, finds the official test directory, runs every test image, writes
the required two-column CSV, checks unique identifiers and nonempty encodings,
and decodes a sample RLE to verify its dimensions. It performs no training,
downloads no model, and calls no inference service.

\section{Evaluation and model selection}\label{sec:evaluation}
Table~\ref{tab:results} lists official public scores displayed for our Kaggle
submissions. They are platform results, not synthetic local estimates. The
0.48 classical run remains the strongest verified route. Both later learned
routes scored lower, so we froze them instead of selecting a more complex
method on architectural appeal alone.

\begin{table}[h]
\caption{Official public leaderboard results for our submissions.}
\label{tab:results}
\centering
\begin{tabular}{lrr}
\toprule
Method & Reference & Score \\
\midrule
Radial normalization + components & 55113085 & \textbf{0.48} \\
Torchvision Mask R-CNN & 55354276 & 0.31 \\
Public YOLO + U-Net refiner & 55354612 & 0.34 \\
\bottomrule
\end{tabular}
\end{table}

The upstream public YOLO notebook reported a score of 0.69 for its original
author. We do not attribute that score to our entry. Our adapted inference run
is the 0.34 result shown above. This distinction is important because the final
submission should be judged from reproducible entrant-owned outputs rather
than a third party's leaderboard record.

The classical method runs in approximately linear time in the number of image
pixels, apart from component bookkeeping, and uses CPU memory for one image and
one integer label map at a time. Unit tests exercise synthetic filament
recovery, disjoint-mask handling, connected-component scoring, and exact COCO
RLE round trips. These tests establish implementation invariants; they are not
substitutes for the official hidden evaluation.

\section{Limitations and reproducibility}
Intensity connectivity is not a semantic model. A filament whose contrast is
interrupted may fragment into several components, while neighboring filaments
that touch after thresholding may merge. Bright active regions and acquisition
artifacts can distort the radial profile. A fixed geometric disk can be less
accurate than per-image limb estimation, and the minimum-area threshold can
discard true fine structures. Finally, public leaderboard performance may not
equal the final evaluation result.

The public release package comprises the MIT-licensed source, exact
\texttt{requirements.txt}, and
\texttt{kaggle/final\_classical\_pipeline.ipynb}. A reviewer attaches the
official competition dataset and runs the notebook from the repository root;
it produces \texttt{/kaggle/working/submission-classical.csv} without any
checkpoint or untracked input. Competition images and generated submissions
are intentionally not redistributed in Git.

AI assistance was used to help implement, test, and document repository
materials. No AI-generated number is presented as an evaluation result. The
entrant remains responsible for executing the notebook, verifying the three
platform references, understanding the method, and submitting the final
repository and report.

\balance
\section{Acknowledgment}
This challenge utilizes the MAGFiLO dataset for evaluation of segmentation
algorithms. We thank the competition organizers, the National Solar
Observatory, and the GONG program for providing the task and observations.

\bibliographystyle{ACM-Reference-Format}
\bibliography{main}
\end{document}
